\begin{afterword}
\summaryhead

The post opens with two tests for the reader: a wheel of fortune stops at 65 before a question about the share of UN countries in Africa, and a product of the numbers 1 to 8 must be estimated in five seconds. Tversky and Kahneman found that subjects who saw 65 gave a median estimate of 45\%, and those who saw 10 gave 25\%. Students shown 1 × 2 × \ldots\ × 8 estimated 512, and those shown the same numbers in descending order estimated 2,250; the true value is 40,320.

The post calls the explanation the current theory: people start from the anchor and adjust until an answer sounds plausible, and so stop too soon. Payoffs for accuracy did not reduce the effect, and absurd anchors such as the year 1215 for Einstein's first visit to America worked as well as plausible ones. Debiasing has mostly failed, the author says, and suggests throwing away an implausible anchor and thinking briefly of one on the other side.

\responsehead

The post reports its studies accurately. The medians for the wheel and for the multiplication task match the 1974 paper, the Einstein anchors match Strack and Mussweiler's table, and anchoring itself has since held up in a large replication. The problem is the explanation the post gives for its main example, and the advice that follows from it.

The post calls adjustment until an answer ``sounds plausible'' the current theory ``for this and similar experiments.'' For anchors that people produce themselves, as in the multiplication task, that was a defended account. For numbers supplied by the experimenter, as with the wheel, the research the post draws on had moved away from it. Strack and Mussweiler, whose Einstein study the post cites, argued that anchoring is a kind of semantic priming, and Epley and Gilovich reported that anchors given by the experimenter appear to work by making anchor-consistent information easier to recall. A month later, in ``Priming and Contamination,'' the author said the same: most anchoring is contamination, not adjustment. Nor does the 1974 wheel study show that people start from the anchor on their own: its subjects were told to move up or down from the given number.

The advice is written mostly for the adjustment account. Readers are told to notice when they are ``adjusting'' a figure and to avoid ``sliding from the anchor,'' and then to dwell briefly on an anchor on the opposite side. The post reports no test of either remedy. A remedy with published support at the time was different: listing reasons why the given anchor is inappropriate, which Mussweiler, Strack and Pfeiffer derived from the accessibility account and found to reduce the effect.

In short: accurate figures for a real effect, explained by a theory that the post's own source had already questioned for its main example, with remedies the post does not test.
\end{afterword}
